% Diagramma di flusso del ciclo di esplorazione (Sezione 1.6)
\begin{tikzpicture}[
    node distance=4.5mm,
    box/.style={rectangle, draw, align=center, minimum width=50mm, minimum height=7mm, font=\small},
    side/.style={rectangle, draw, align=center, minimum width=28mm, minimum height=7mm, font=\small},
    term/.style={rectangle, draw, rounded corners=3.5mm, align=center, minimum width=28mm, minimum height=7mm, font=\small},
    dec/.style={diamond, draw, aspect=2.6, align=center, inner sep=0.5pt, font=\small},
    arr/.style={-{Stealth[length=2mm]}},
    lab/.style={font=\footnotesize}
]
    % colonna principale
    \node[term] (start) {Inizio};
    \node[box, below=of start] (scan) {Scansione e aggiornamento della mappa};
    \node[dec, below=of scan] (path) {Percorso\\da seguire?};
    \node[dec, below=of path] (front) {Esistono\\frontiere?};
    \node[dec, below=of front] (reach) {Frontiere\\raggiungibili?};
    \node[box, below=of reach] (choose) {Scelta della destinazione\\(strategia di esplorazione)};
    \node[box, below=of choose] (plan) {Calcolo del percorso (Dijkstra)};
    \node[dec, below=of plan] (next) {Cella successiva\\libera?};
    \node[box, below=of next] (move) {Spostamento di una cella};

    % esiti laterali
    \node[term] (done) at ($(front)+(58mm,0)$) {Esplorazione\\completa};
    \node[term] (block) at ($(reach)+(58mm,0)$) {Arresto\\(blocco)};
    \node[side] (drop) at ($(next)+(58mm,0)$) {Abbandono\\del percorso};

    % linee verticali dei ritorni
    \coordinate (L1) at (-37mm,0);
    \coordinate (L2) at (-44mm,0);
    \coordinate (R) at (80mm,0);
    \coordinate (J) at ($(path.south)!0.65!(front.north)$);

    % flusso principale
    \draw[arr] (start) -- (scan);
    \draw[arr] (scan) -- (path);
    \draw[arr] (path) -- node[lab, left, pos=0.3] {no} (front);
    \draw[arr] (front) -- node[lab, right] {sì} (reach);
    \draw[arr] (reach) -- node[lab, right] {sì} (choose);
    \draw[arr] (choose) -- (plan);
    \draw[arr] (plan) -- (next);
    \draw[arr] (next) -- node[lab, right] {sì} (move);

    % uscite laterali
    \draw[arr] (front) -- node[lab, above] {no} (done);
    \draw[arr] (reach) -- node[lab, above] {no} (block);
    \draw[arr] (next) -- node[lab, above] {no} (drop);

    % percorso gia' disponibile: si passa direttamente al controllo della cella successiva
    \draw[arr] (path.west) -- node[lab, above, pos=0.25] {sì} (path.west -| L1) |- (next.west);
    % dopo lo spostamento si esegue una nuova scansione
    \draw[arr] (move.west) -- (move.west -| L2) |- (scan.west);
    % dopo l'abbandono del percorso si torna alla scelta della destinazione
    \draw[arr] (drop.east) -- (drop.east -| R) |- (J);
\end{tikzpicture}
